\paragraph{Sentence diversity and the limits on answer length.}
We compared four passes over 48,000 fresh training pairs with a single pass over 192,000 distinct pairs. The larger corpus raised the number of exact translations by one on each 2,048-request cohort, from 31 to 32 and from 19 to 20. Across the two cohorts, six of its 52 exact answers required changing the source. Aggregate glyph error increased by 0.21 and 0.19 percentage points (paired chapter-bootstrap 95\% intervals for the increases, 0.03--0.43 and 0.07--0.32). Both arms copied all 4,085 renderable quotations exactly. With correct previous actions and writing, the larger-corpus editor drew 11.25\% and 11.01\% of required replacement or insertion bitmaps exactly. Given the target ink measurements, its renderer recovered 99.96\% and 99.95\% (Fig.~\ref{fig:edit-diversity}). The model can write a legible sentence whose meaning differs from its source.

The frozen policy limited the number of insertions at each source position to the most probable learned count. Even with perfect glyph choices, this limit made 1,913 of 2,043 and 1,932 of 2,042 reference lengths unreachable. At least 44,211 of the observed 76,869 glyph edits were unavoidable under this length constraint. No reference in these cohorts exceeded the global 192-cell ceiling. Correcting the insertion decision and learning what to write are therefore distinct empirical questions.

We then changed only the permitted insertion count, taking the conditional 90th percentile at each cursor. This reduced unreachable reference lengths to 51 and 43. Despite the increased capacity, discovery and replication error rose from 56.90\% and 55.97\% to 68.61\% and 67.87\%. Exact translations fell from 32 and 20 to 23 and 18. Increasing the permitted length allowed the same weights to emit longer answers with more incorrect writing. This follow-up used the same cohorts that motivated the diagnosis.
